\documentclass[11pt]{article}

% ============================================================
% Packages
% ============================================================
\usepackage[utf8]{inputenc}
\usepackage[margin=1in]{geometry}
\usepackage{amsmath,amssymb,amsthm}
\usepackage{graphicx}
\usepackage{booktabs}
\usepackage{xcolor}
\usepackage{hyperref}
\usepackage{tikz}
\usepackage{enumitem}
\usepackage{pgfplots}

\pgfplotsset{compat=1.18}

% ============================================================
% Document Formatting
% ============================================================
\setlength{\parskip}{0.6em}
\setlength{\parindent}{0pt}

% ============================================================
% List Formatting
% ============================================================
\setlist[enumerate]{
    leftmargin=2em,
    itemsep=0.25em
}

% ============================================================
% Solution Environment
% ============================================================
\newenvironment{solution}
    {\par\noindent\textbf{Solution.}\par}
    {\par}

% ============================================================
% Link Styling
% ============================================================
\hypersetup{
    colorlinks=true,
    linkcolor=blue,
    urlcolor=cyan
}

% ============================================================
% Metadata
% ============================================================
\title{Ordinary Differential Equations}
\author{Chris Ritter}
\date{September 8, 2026}

% ============================================================
% Document
% ============================================================

\begin{document}

\maketitle
\tableofcontents
\newpage

% ============================================================
\section{Linear and Non-Linear Equations}
% ============================================================

For a linear initial-value problem (IVP),
\[
\frac{dy}{dt}=-p(t)y+g(t), \qquad y(t_0)=y_0.
\]
More generally, an IVP may have the form
\[
\frac{dy}{dt}=f(t,y), \qquad y(t_0)=y_0.
\]

Two questions arise:
\begin{enumerate}
    \item Does the IVP have at least one solution (existence)?
    \item If so, is that solution unique?
\end{enumerate}

\subsection{Existence and Uniqueness for Linear Equations}

If \(p\) and \(g\) are continuous on an interval
\(I=(\alpha,\beta)\) containing \(t_0\), then the linear IVP has a
unique solution defined throughout \(I\).

With
\[
\mu(t)=\exp\!\left(\int_{t_0}^{t}p(s)\,ds\right),
\]
the solution can be written as
\[
y(t)=\frac{y_0+\displaystyle\int_{t_0}^{t}\mu(s)g(s)\,ds}{\mu(t)}.
\]

For uniqueness, suppose \(u\) and \(v\) satisfy the same equation and
initial condition. Their difference \(w=u-v\) satisfies
\[
w'=-p(t)w, \qquad w(t_0)=0.
\]
Multiplying by \(\mu(t)\) gives
\[
\frac{d}{dt}\bigl(\mu(t)w(t)\bigr)=0.
\]
Consequently, \(\mu(t)w(t)=\mu(t_0)w(t_0)=0\). Since
\(\mu(t)>0\) throughout \(I\), we have \(w=0\), so \(u=v\).
The formula above also proves existence: differentiating it recovers
the differential equation and setting \(t=t_0\) gives \(y(t_0)=y_0\).

\subsection{Existence and Uniqueness for Nonlinear Equations}

For \(y'=f(t,y)\), continuity of \(f\) near the initial point
guarantees a solution at least locally. If \(f\) and
\(\partial f/\partial y\) are continuous on a rectangle around that
point, the local solution is unique. These are local statements: a
solution may cease to exist at a finite endpoint of its maximal
interval.

\subsection{Example: A Linear IVP}

Does \(ty'+2y=4t^2\), \(y(1)=2\), have a unique solution?

\begin{solution}
For \(t\ne0\), divide by \(t\):
\[
y'+\frac{2}{t}y=4t.
\]
On \(I=(0,\infty)\), both coefficients are continuous. The
integrating factor is \(\mu(t)=t^2\), so
\[
(t^2y)'=4t^3, \qquad t^2y=t^4+C.
\]
Thus \(y=t^2+C/t^2\). The initial condition gives \(C=1\), hence
\[
\boxed{y(t)=t^2+\frac{1}{t^2},\qquad 0<t<\infty.}
\]
The linear theorem guarantees uniqueness on \((0,\infty)\). This
solution cannot be extended through \(t=0\), where it diverges.
\end{solution}

\subsection{Example: A Nonlinear IVP with a Singular Boundary}

Consider
\[
\frac{dy}{dx}=\frac{3x^2+4x+2}{2(y-1)}, \qquad y(0)=-1.
\]

\begin{solution}
The right-hand side and its derivative with respect to \(y\) are
continuous near \((0,-1)\), since \(y\ne1\) there. A unique local
solution therefore exists. Separating and integrating yields
\[
2(y-1)\,dy=(3x^2+4x+2)\,dx,
\qquad
(y-1)^2=x^3+2x^2+2x+C.
\]
Since \(y(0)=-1\), we have \(C=4\). Taking the negative square-root
branch to match the initial value gives
\[
\boxed{y(x)=1-\sqrt{(x+2)(x^2+2)},\qquad -2<x<\infty.}
\]
At \(x=-2\), this formula reaches \(y=1\), where the original
equation is undefined. Thus \((-2,\infty)\) is the maximal interval
containing \(0\) on which this solution satisfies the IVP.
\end{solution}

\subsection{Example: Existence Without Uniqueness}

Consider
\[
\frac{dy}{dt}=y^{1/3}, \qquad y(0)=0, \quad t\geq0.
\]

\begin{solution}
The function \(f(y)=y^{1/3}\) is continuous everywhere, so a local
solution exists. However, \(f'(y)=\frac13y^{-2/3}\) is undefined at
\(y=0\), so the usual uniqueness criterion does not apply.
Indeed, both
\[
y(t)=0
\qquad\text{and}\qquad
y(t)=\left(\frac{2t}{3}\right)^{3/2}
\]
satisfy the IVP for \(t\geq0\). More generally, for any \(a\geq0\),
\[
y_a(t)=
\begin{cases}
0, & 0\leq t\leq a,\\
\left(\dfrac{2(t-a)}{3}\right)^{3/2}, & t>a
\end{cases}
\]
is a solution. The derivative from the right at \(t=a\) is zero, so
these pieces join differentiably. There are infinitely many solutions.
\end{solution}

\subsection{Example: Unique Locally, but Not for All Time}

Consider
\[
\frac{dy}{dx}=y^2, \qquad y(0)=1.
\]

\begin{solution}
Both \(f(x,y)=y^2\) and \(\partial f/\partial y=2y\) are continuous
everywhere, so the IVP has a unique local solution. Separating gives
\[
-\frac{1}{y}=x+C.
\]
Applying \(y(0)=1\) gives
\[
\boxed{y(x)=\frac{1}{1-x},\qquad -\infty<x<1.}
\]
The solution grows without bound as \(x\to1^-\). Its maximal
interval containing the initial point is \((-\infty,1)\), even though
the equation's right-hand side is defined for every \((x,y)\).
\end{solution}

\end{document}
